\documentclass[11pt]{article}

\usepackage[margin=1in]{geometry}
\usepackage[T1]{fontenc}
\usepackage[utf8]{inputenc}
\usepackage{mathpazo}
\usepackage{microtype}
\usepackage{amsmath,amssymb,amsthm,mathtools}
\usepackage{booktabs,tabularx,array}
\usepackage{enumitem}
\usepackage{xcolor}
\usepackage{hyperref}
\usepackage{cleveref}

\definecolor{linkblue}{HTML}{163A5F}
\hypersetup{
  colorlinks=true,
  linkcolor=linkblue,
  citecolor=linkblue,
  urlcolor=linkblue,
  pdftitle={A Minimal Two-Firm Example of Causal Dissonance and Awareness Refinement},
  pdfauthor={Research note}
}

\newtheorem{definition}{Definition}
\newtheorem{assumption}{Assumption}
\newtheorem{proposition}{Proposition}
\newtheorem{remark}{Remark}

\newcommand{\ind}{\mathbf 1}
\newcommand{\emptysetset}{\varnothing}

\title{A Minimal Two-Firm Example of\\
Causal Dissonance and Awareness Refinement}
\author{Research note}
\date{July 2026}

\begin{document}

\maketitle

\begin{abstract}
This note develops a deliberately minimal deterministic example in which heterogeneous causal beliefs generate heterogeneous strategic actions; those actions expose a contradiction within the agents' currently expressible model class; and inquiry resolves the contradiction by refining an awareness partition.  The refinement makes previously unthinkable causal theories expressible, after which Bayesian judgment can resume.  The note is a formal building block rather than a paper model: it does not yet explain how inquiry selects among several possible refinements, how a new prior is derived from an old one, or how inquiry incentives interact strategically.
\end{abstract}

\section{Purpose and modeling choices}

The example separates three objects that should not be conflated:
\begin{enumerate}[label=(\roman*)]
  \item an \emph{awareness partition}, which determines the customer distinctions an agent can represent;
  \item a \emph{causal theory}, which assigns customer responses to the cells of that partition; and
  \item a \emph{strategy}, which maps represented customer states into product choices.
\end{enumerate}
No separate formal object for language is introduced.  The partition is interpreted as being induced by the predicates available to the agent when formulating causal theories.  Likewise, the question generated by a contradiction is given a verbal interpretation rather than separate notation.  In the two-state case the relevant partition refinement is unique, so no general inquiry-selection technology is needed.

\section{Notation}

\begin{center}
\renewcommand{\arraystretch}{1.18}
\begin{tabularx}{\textwidth}{@{}>{\raggedright\arraybackslash}p{0.25\textwidth}X@{}}
\toprule
\textbf{Symbol} & \textbf{Meaning} \\
\midrule
$i\in\{A,B\}$ & Firm index. \\
$X=\{x_1,x_2\}$ & Finite customer-state space. \\
$x_t\in X$ & Customer state in period $t$. \\
$A=\{0,1\}$ & Set of product architectures and firm actions. \\
$a_{it}\in A$ & Architecture offered by firm $i$ in period $t$. \\
$d:X\to\{0,1\}$ & Objectively preferred architecture at each customer state. \\
$u_i$ & Firm $i$'s deterministic objective payoff function. \\
$u_{it}$ & Firm $i$'s realized payoff in period $t$. \\
$\mathcal P_i$ & Firm $i$'s awareness partition of $X$. \\
$M(\mathcal P_i)$ & Set of causal theories expressible under $\mathcal P_i$. \\
$m\in M(\mathcal P_i)$ & An individual causal theory. \\
$\mu_i\in\Delta(M(\mathcal P_i))$ & Firm $i$'s belief over its expressible theories. \\
$\sigma_i:\mathcal P_i\to A$ & Firm $i$'s pure strategy. \\
$h_{it}$ & Firm $i$'s observed history through period $t$. \\
$\widehat M_{it}$ & Active theories consistent with $h_{it}$. \\
$\tau_i$ & First period in which $\widehat M_{it}=\emptysetset$. \\
\bottomrule
\end{tabularx}
\end{center}

Capital Roman letters denote sets, lowercase Roman and Greek letters denote elements, parameters, and functions, calligraphic letters denote sets of sets, and blackboard-bold notation is reserved for canonical sets.  For any finite set $Z$, $\Delta(Z)$ denotes the set of probability distributions on $Z$.

\section{Objective environment}

There are two firms, $A$ and $B$, and two customer states,
\[
X=\{x_1,x_2\}.
\]
The two states are analyst-level descriptions.  They represent customer situations for which different product architectures create greater value.

The firms have a common action set
\[
A=\{0,1\}.
\]
Choosing action $a_i$ means offering architecture $a_i$; architecture is not a separate object from action.

The objectively preferred architecture is described by
\[
d:X\longrightarrow\{0,1\},
\qquad
d(x_1)=0,
\qquad
d(x_2)=1.
\]
Thus architecture 0 better serves customer state $x_1$, whereas architecture 1 better serves customer state $x_2$.

The customer has no outside option.  If the firms offer different architectures, the customer buys from the firm offering $d(x)$.  If the firms offer the same architecture, they split the market, even when a different architecture would have served the customer better.  Firm $A$'s payoff is
\begin{equation}
u_A(a_A,a_B,x)=
\begin{cases}
\frac12, & a_A=a_B,\\
1, & a_A\neq a_B\text{ and }a_A=d(x),\\
0, & a_A\neq a_B\text{ and }a_B=d(x).
\end{cases}
\label{eq:objective-payoff}
\end{equation}
Firm $B$'s payoff is
\[
u_B(a_A,a_B,x)=1-u_A(a_A,a_B,x).
\]
The realized period-$t$ payoff is
\[
u_{it}=u_i(a_{At},a_{Bt},x_t).
\]
All causal relationships are deterministic.  No arrival distribution or outcome disturbance is imposed.  The example simply studies the two-period sequence $x_1,x_2$.

\section{Awareness partitions and expressible theories}

\begin{definition}[Awareness partition]
An awareness partition $\mathcal P_i$ is a partition of $X$.  Customer states in the same cell are not distinguishable within firm $i$'s current causal representation.  A pure strategy is a function
\[
\sigma_i:\mathcal P_i\longrightarrow A.
\]
Consequently, the strategy must prescribe the same action at every customer state in a given cell.
\end{definition}

Initially both firms have the coarse partition
\[
\mathcal P_i^0=\{X\}=\big\{\{x_1,x_2\}\big\}.
\]
The analyst distinguishes $x_1$ from $x_2$, but the firms' initial causal representations do not.  A feasible initial pure strategy is therefore constant on $X$.

Let $n$ denote the claim that architecture is causally neutral.  For any partition $\mathcal P$ of $X$, define the expressible model set
\begin{equation}
M(\mathcal P)
=
\left\{
m:X\longrightarrow\{0,1,n\}:
m(x)=m(x')
\text{ whenever }x,x'\text{ lie in the same cell of }\mathcal P
\right\}.
\label{eq:model-space}
\end{equation}
The interpretations are
\[
m(x)=0:
\text{ architecture 0 is preferred at }x,
\]
\[
m(x)=1:
\text{ architecture 1 is preferred at }x,
\]
and
\[
m(x)=n:
\text{ architecture is irrelevant at }x.
\]

The payoff predicted by theory $m$ is
\begin{equation}
u_A^m(a_A,a_B,x)=
\begin{cases}
\frac12, & a_A=a_B\text{ or }m(x)=n,\\
1, & a_A\neq a_B\text{ and }m(x)=a_A,\\
0, & a_A\neq a_B\text{ and }m(x)=a_B,
\end{cases}
\label{eq:subjective-payoff}
\end{equation}
with $u_B^m=1-u_A^m$.

Under the initial partition, the only expressible theories are the three constant functions
\[
M(\mathcal P_i^0)=\{m^0,m^1,m^n\},
\]
where, for every $x\in X$,
\[
m^0(x)=0,
\qquad
m^1(x)=1,
\qquad
m^n(x)=n.
\]
The true theory $d$ is not an element of $M(\mathcal P_i^0)$ because it is not constant on $X$.

\section{Private priors and simultaneous product choice}

Firm $i$ has a private prior
\[
\mu_i^0\in\Delta\big(M(\mathcal P_i^0)\big).
\]
The priors need not be common knowledge, and neither firm requires a theory of why its rival holds a different prior.

\begin{assumption}[Heterogeneous causal beliefs]
\[
\mu_A^0(m^0)>\mu_A^0(m^1)
\qquad\text{and}\qquad
\mu_B^0(m^1)>\mu_B^0(m^0).
\]
\end{assumption}

In each period the firms choose architectures simultaneously.  Neither observes the rival's current action before choosing.  For either firm $i$, and against either rival action, the subjective expected-payoff difference between architectures 0 and 1 is
\begin{equation}
\frac12\left[\mu_i^0(m^0)-\mu_i^0(m^1)\right].
\label{eq:dominance-difference}
\end{equation}
The mass on $m^n$ cancels because that model assigns payoff $1/2$ to either architecture.

\begin{proposition}[Belief-induced differentiation]
Under heterogeneous causal beliefs, architecture 0 is subjectively strictly dominant for firm $A$, architecture 1 is subjectively strictly dominant for firm $B$, and the simultaneous action profile in both periods is
\[
(a_{At},a_{Bt})=(0,1).
\]
\end{proposition}

\begin{proof}
Equation~\eqref{eq:dominance-difference} is positive for firm $A$ and negative for firm $B$.  Each conclusion is independent of the rival's current action.
\end{proof}

\section{Experience, active theories, and causal dissonance}

Firm $i$'s observed history through period $t$ is
\[
h_{it}=\big(x_s,a_{As},a_{Bs},u_{is}\big)_{s=1}^t.
\]
Although a customer state is stored as part of the history, a theory in $M(\mathcal P_i^0)$ cannot assign different causal significance to $x_1$ and $x_2$.

Define the active model set
\begin{equation}
\widehat M_{it}
=
\left\{
m\in M(\mathcal P_i^0):
u_i^m(a_{As},a_{Bs},x_s)=u_{is}
\text{ for every }s\leq t
\right\}.
\label{eq:active-set}
\end{equation}
The expressible model set $M(\mathcal P_i^0)$ never becomes empty.  Rather, $\widehat M_{it}$ records which currently expressible theories remain consistent with the observed history.

Whenever $\widehat M_{it}\neq\emptysetset$, deterministic Bayesian updating gives
\begin{equation}
\mu_{it}(m)
=
\frac{
\mu_i^0(m)\ind\{m\in\widehat M_{it}\}
}{
\sum_{m'\in M(\mathcal P_i^0)}
\mu_i^0(m')\ind\{m'\in\widehat M_{it}\}
}.
\label{eq:bayes-update}
\end{equation}

The realized two-period history is shown in \cref{tab:history}.

\begin{table}[ht]
\centering
\caption{The two-period deterministic history}
\label{tab:history}
\renewcommand{\arraystretch}{1.15}
\begin{tabular}{@{}ccccccc@{}}
\toprule
Period & State & $d(x_t)$ & $a_{At}$ & $a_{Bt}$ & $u_{At}$ & $u_{Bt}$ \\
\midrule
1 & $x_1$ & 0 & 0 & 1 & 1 & 0 \\
2 & $x_2$ & 1 & 0 & 1 & 0 & 1 \\
\bottomrule
\end{tabular}
\end{table}

After period 1,
\[
\widehat M_{A1}=\widehat M_{B1}=\{m^0\}.
\]
Both firms therefore become certain of $m^0$ within their current model spaces.  Period 2 then falsifies $m^0$, yielding
\[
\widehat M_{A2}=\widehat M_{B2}=\emptysetset.
\]
The Bayesian posterior in \eqref{eq:bayes-update} is undefined because its denominator is zero.

Define the gap time
\[
\tau_i=\min\{t:\widehat M_{it}=\emptysetset\}.
\]
In the example,
\[
\tau_A=\tau_B=2.
\]

\begin{definition}[Causal dissonance]
Periods $s<t$ generate causal dissonance for firm $i$ when
\begin{enumerate}[label=(\alph*)]
  \item $x_s$ and $x_t$ lie in the same cell of $\mathcal P_i$;
  \item $(a_{As},a_{Bs})=(a_{At},a_{Bt})$; and
  \item $u_{is}\neq u_{it}$.
\end{enumerate}
\end{definition}

Causal dissonance is evidence that the firm's represented causal antecedents are insufficient for a deterministic account of the outcome.  It is not itself evidence identifying the omitted cause.  The immediate interpretive question is: \emph{what currently unrepresented difference between the customer states explains the payoff reversal?}

\begin{proposition}[Differentiation and dissonance]
Suppose both customer states occur.  If the firms choose different architectures at both states, each firm experiences causal dissonance and its initial active model set becomes empty.  If the firms choose the same architecture at both states, both receive $1/2$ at both states; causal dissonance does not arise, and all three initial theories remain active.
\end{proposition}

\begin{proof}
When the actions differ, the firm offering architecture 0 wins at $x_1$ and loses at $x_2$; the other firm has the reverse experience.  No constant deterministic theory fits both outcomes.  When the actions coincide, \eqref{eq:objective-payoff} gives each firm $1/2$ at both states, which every theory in $M(\mathcal P_i^0)$ predicts under coincident actions.
\end{proof}

The second case is a self-confirming awareness trap: the firms receive no indication that their coarse representations prevent a profitable state-contingent strategy.

\section{Inquiry as restorative partition refinement}

For two partitions $\mathcal P$ and $\mathcal P'$ of $X$, write
\[
\mathcal P\prec\mathcal P'
\]
when $\mathcal P'$ strictly refines $\mathcal P$: every cell of $\mathcal P'$ is contained in a cell of $\mathcal P$, and the partitions are not equal.

For any proposed partition $\mathcal P'$, define the models that would be active under that partition as
\begin{equation}
\widehat M_i(h_{it};\mathcal P')
=
\left\{
m\in M(\mathcal P'):
u_i^m(a_{As},a_{Bs},x_s)=u_{is}
\text{ for every }s\leq t
\right\}.
\label{eq:counterfactual-active-set}
\end{equation}

\begin{definition}[Restorative refinement]
A partition $\mathcal P'$ is a restorative refinement of $\mathcal P_i$ at history $h_{it}$ if
\[
\mathcal P_i\prec\mathcal P'
\]
and
\[
\widehat M_i(h_{it};\mathcal P')\neq\emptysetset.
\]
A restorative refinement is coarsest if no strictly intermediate partition is also restorative.
\end{definition}

The restoration condition does the substantive work: the new distinctions must permit a stable causal theory that accounts for the dissonant experience.  A constant theory or a refinement that leaves the two dissonant customer states in the same cell cannot satisfy this condition.  Because a partition of $X$ defines distinctions available before action, prospective applicability is built into the representation.

In the two-state example there is exactly one strict refinement of $\mathcal P_i^0$:
\begin{equation}
\mathcal P_i^1
=
\big\{\{x_1\},\{x_2\}\big\}.
\label{eq:refined-partition}
\end{equation}
It is restorative and therefore uniquely coarsest.  Conditional on successful inquiry, the insight is the awareness transition
\[
\mathcal P_i^0
=
\big\{\{x_1,x_2\}\big\}
\quad\longrightarrow\quad
\mathcal P_i^1
=
\big\{\{x_1\},\{x_2\}\big\}.
\]
No separate question space, candidate-classifier space, or partition-selection rule is required in this example.

\section{New theories and Bayesian judgment}

Under the refined partition,
\[
M(\mathcal P_i^1)
=
\{m:X\to\{0,1,n\}\}.
\]
There are $3^2=9$ expressible theories.  In particular,
\[
d\in M(\mathcal P_i^1)
\qquad\text{but}\qquad
d\notin M(\mathcal P_i^0).
\]
The awareness refinement therefore makes the objectively correct theory expressible without treating it as an element of the initial model space.

After refinement, firm $i$ forms a prior
\[
\mu_i^1\in\Delta\big(M(\mathcal P_i^1)\big).
\]
Ordinary Bayesian updating does not determine $\mu_i^1$ from $\mu_i^0$ because the domain of belief has changed.  For the deterministic example, assume only that $\mu_i^1$ has full support.

Reconsidering the two-period history yields
\[
\widehat M_i(h_{i2};\mathcal P_i^1)=\{d\}.
\]
The first observation requires $m(x_1)=0$, while the second requires $m(x_2)=1$.  Thus Bayesian judgment on the expanded model space gives
\[
\mu_{i2}^1(d)=1.
\]
The rapid convergence is specific to the deterministic two-state example.  With additional states, limited histories, or stochastic outcomes, several refined theories can remain active and the selected refinement can be wrong.

\begin{proposition}[Unique restorative insight and truth recovery]
In the two-state deterministic example, causal dissonance under differentiated actions has a unique restorative partition refinement.  If the post-refinement prior has full support, Bayesian reconsideration of the observed history identifies the objectively correct theory $d$.
\end{proposition}

\begin{proof}
The singleton partition in \eqref{eq:refined-partition} is the only strict refinement of $\mathcal P_i^0$.  It permits state-contingent theories.  The two differentiated-action outcomes uniquely require $m(x_1)=0$ and $m(x_2)=1$, so the active set contains only $d$.
\end{proof}

\section{Strategic value of asymmetric insight}

Suppose firm $A$ refines and learns $d$ before firm $B$, while firm $B$ continues to offer architecture 1 at both customer states.  Firm $A$ can use the contingent strategy
\[
\sigma_A(\{x\})=d(x).
\]
It wins at $x_1$ and splits the market at $x_2$.  Its average payoff across the two states is therefore
\[
\frac12(1)+\frac12\left(\frac12\right)=\frac34,
\]
while firm $B$ receives $1/4$.  Before the refinement, differentiated constant strategies give each firm an average payoff of $1/2$.  A uniquely aware firm therefore earns an insight rent of $1/4$ in this example.

If both firms refine, learn $d$, and offer the preferred architecture at each state, they split both markets and each returns to an average payoff of $1/2$.  The insight rent is transient unless awareness remains asymmetric.

\section{What the example establishes---and what it does not}

The example establishes a coherent sequence:
\begin{align*}
\text{private causal beliefs}
&\longrightarrow
\text{strategic differentiation}
\longrightarrow
\text{causal dissonance}\\
&\longrightarrow
\text{restorative refinement}
\longrightarrow
\text{expanded theory space}\\
&\longrightarrow
\text{Bayesian judgment}.
\end{align*}
It also establishes a strategically meaningful contrast: differentiated actions reveal the inadequacy of coarse awareness, whereas identical actions can sustain a self-confirming awareness trap.

The example deliberately does not yet model:
\begin{enumerate}[label=(\roman*)]
  \item how inquiry selects among several coarsest restorative refinements;
  \item cognitive or organizational restrictions on reachable refinements;
  \item the cost, timing, or expected benefit of inquiry;
  \item a principled relationship between $\mu_i^0$ and $\mu_i^1$;
  \item noisy outcomes or stochastic causal laws;
  \item strategic concealment, induced dissonance, or rival learning; or
  \item whether a refinement that rationalizes past experience generalizes correctly.
\end{enumerate}

Those omissions identify the next theoretical problem rather than defects in the minimal example.  Once $X$ contains more than two states, several restorative refinements may fit the same history.  A theory of directed inquiry must then explain which refinement becomes available, why that refinement is selected, and how the agent forms beliefs over the newly expressible models without treating them as elements of its original awareness space.

\section{Conclusion}

The minimal example's main contribution is architectural.  An awareness partition determines the theories an agent can formulate.  Experience can empty the active subset of those theories without directly supplying a replacement.  Inquiry responds by refining the awareness partition; the refinement expands the expressible model space; and only then can Bayesian judgment operate over the new possibilities.  Strategic interaction is essential because the agents' causal beliefs determine whether their actions generate revealing dissonance or suppress it.

\end{document}
